\documentclass[12pt]{article}
\usepackage[utf8]{inputenc}
\usepackage[T1]{fontenc}
\usepackage{hyperref}
\title{Neve'ei Language Resources for K-12 Curriculum Development}
\author{Compiled by Claude Sonnet 5 (claude-sonnet-5) for Andr\'as Kornai}
\date{}
\begin{document}
\maketitle

\section{Introduction}
Neve'ei (ISO 639-3 vnm; also Vinmavis, Lambumbu) is an Oceanic Austronesian language of central Malekula, Vanuatu, with about 500 speakers. It has no standard orthography; the spelling in Musgrave's grammar is the most widely referenced.

\section{Primary Educational Resources}
No primers, textbooks or school programmes were found.

\section{Linguistic and Reference Materials}
\begin{itemize}
\item Jill Musgrave, \emph{A grammar of Neve'ei, Vanuatu} (Pacific Linguistics 587, 2007; open access, CC BY-SA 4.0): \url{https://openresearch-repository.anu.edu.au/items/54461ea4-c24c-43d8-8d7d-094986566440/full}
\item PDF of the same volume: \url{https://openresearch-repository.anu.edu.au/bitstream/1885/146743/1/587_Musgrave.pdf}
\item Austronesian Basic Vocabulary Database wordlist; it cites an unpublished Crowley and Musgrave Neve'ei dictionary manuscript that is not publicly available: \url{https://abvd.eva.mpg.de/austronesian/language.php?id=484}
\item Wikipedia article: \url{https://en.wikipedia.org/wiki/Neve%CA%BBei_language}
\item Verbix short verb list (not a dictionary): \url{https://docs.verbix.com/Languages/Neveei}
\end{itemize}

\section{Religious and Community Resources}
None were found.

\section{Digital and Audio Resources}
\begin{itemize}
\item Endangered Languages Project entry: \url{https://endangeredlanguages.com/node/97182}
\item Omniglot page on the writing and numbers: \url{https://omniglot.com/writing/neveei.htm}
\end{itemize}

\section{Curriculum Development Context}
Most adults speak the language, but not all young people do, and no school use is known. The open-access grammar and the data in it give a starting point for an orthography and teaching material. No published dictionary or reading materials were found, so these would have to be created with community involvement.
\end{document}